\documentclass[10pt,a4paper]{article}
\usepackage[utf8]{inputenc}
\usepackage{booktabs}
\usepackage[a4paper,left=3cm,right=2cm]{geometry}
\usepackage{amsmath}
\usepackage[thmmarks, amsmath, thref]{ntheorem}
\usepackage{amsfonts}
\usepackage{amssymb}
\usepackage{fancyhdr}
\pagestyle{fancy}
\usepackage{anysize}
\usepackage{graphicx}
\usepackage{float}
\usepackage{hyperref}
\usepackage{enumitem}
\usepackage{verbatim}
\newtheorem{theorem}{Theorem}
\theoremsymbol{\ensuremath{\blacksquare}}
\newtheorem*{solution}{Solución:}

\usepackage{xcolor}

% Margenes y formalidades

\textwidth 6.25in % Margen Derecho y ancho del texto
\textheight 8.5in
\oddsidemargin 0in % Margen Izquierdo
\headheight 0in % Largo del texto
\leftmargin 1in
\parindent 0em % Salto entre la indentación (Sangría)https://www.overleaf.com/4519127669sdpptkhvsbfn
\topmargin -0.5in % Margen Superior
\parskip 0.5ex % Salto entre parrafos
\baselineskip 1pt



\lhead{
\includegraphics[scale=0.08]{Figuras/logoUSM.png}} % texto izquierda de la cabecera
%\chead{}                                      % texto centro de la cabecera
\rhead{Universidad Técnica Federico Santa María\\Ingeniería Comercial } % número de página a la derecha
\renewcommand{\headrulewidth}{0.4pt} % grosor de la línea de la cabecera
\renewcommand{\footrulewidth}{0.4pt} % grosor de la línea del pie

\begin{document}
\vspace{5cm}



\begin{center}
\begin{large}
\textbf{\\CERTAMEN 2:}\textbf{ ECONOMETRÍA ICS-294}
\end{large}\\
\vspace{0.2cm}
II Semestre 2025\\
Profesores: Jocelyn Tapia Stefanoni, Franco Antonucci Fernández y Ricardo Pruneda Urz\'ua\\
%Ayudante: 
\end{center}

\textbf{Instrucciones:} 

\begin{itemize}

    \item Usted dispone de 90 minutos  para intentar los 100 puntos del certamen.
    \item Cuide la presentación y redacción de sus preguntas.
    \item No puede utilizar celular, computador o tablet, pero sí una calculadora TI.
     
\end{itemize}

\section{Comentes y preguntas cortas (40 puntos)}

Responda solo 4 de las siguientes preguntas.
\begin{enumerate}
    \item \textbf{RdeA9.1} (10 puntos) 
   Muestre que si se cumplen los supuestos de Teorema de Gauss-Markov (los primeros 4) el estimador de MCO $ \hat{\beta} = (X^{\top}X)^{-1}X^{\top}y$
es insesgado. Señale como utilizó los supuestos en su demostración.

{\color{blue} \begin{align*}
        \hat{\beta} &= (X^{\top}X)^{-1}X^{\top}y \quad \text{linealidad- no multicolinealidad perfecta} \\
        &= (X^{\top}X)^{-1}X^{\top}(X\beta + u) \quad \text{} \\
        &= \beta + (X^{\top}X)^{-1}X^{\top}u \quad  \\
        E(\hat{\beta}|X) &= \beta + (X^{\top}X)^{-1}X^{\top}E(u|X) \quad  \\
        &= \beta \quad \text{muestra aleatoria- ortogonalidad}
    \end{align*}}
 \item \textbf{RdeA9.1} (10 puntos)  
Suponga un modelo dado por:

$$\text{PAES}_i = \beta_0 + \beta_1 \cdot expend_i + \beta_2 \cdot ingreso_i + u$$

Donde PAES es el puntaje de ingreso que un alumno obtuvo en la prueba de acceso a la educación superior, \textit{expend} es el gasto en educación en el estudiante e \textit{ingreso} es el ingreso familiar promedio del estudiante. 
Se esperaría que \textit{expend} e \textit{ingreso} estén correlacionadas, puesto que las comunas donde en los que el ingreso familiar promedio es mayor tienden a gastar más en educación por estudiante. Es por esto que el modelo no cumple con el supuesto de colinealidad y, por lo tanto, las estimaciones serán insesgadas e inconsistentes. Comente.

 
{\color{blue}
Falso. La correlación entre las variables \textit{expend} e \textit{ingreso} (multicolinealidad)  no genera sesgo ni inconsistencia en los estimadores de Mínimos Cuadrados Ordinarios (MCO).
No obstante, la multicolinealidad  sí afecta la \textbf{precisión de las estimaciones}, aumentando las varianzas de los estimadores y reduciendo la significancia estadística de los coeficientes individuales. Es decir, los intervalos de confianza serán más amplios y los estadísticos t menos precisos, dificultando la inferencia.
En el caso extremo de \textbf{multicolinealidad perfecta} (una relación lineal exacta entre \textit{expend} e \textit{ingreso}), el modelo no podría estimarse por MCO, ya que las variables explicativas serían linealmente dependientes.
}


 \item \textbf{RdeA9.1} (10 puntos) 
En la tabla 1  se muestran los resultados de una investigación que busca identificar los determinantes de los salarios de los profesores (avgsal).  Utilizan las siguientes variables explicativas 

\begin{itemize}
   \item \texttt{lunch}: porcentaje de estudiantes elegibles para almuerzo gratuito.
  \item \texttt{enrol}: matrícula (número de estudiantes).
  \item \texttt{staff}: personal por cada 1000 estudiantes.
 
  \item \texttt{math4}: porcentaje que aprueba la prueba de matemáticas de 4.\textsuperscript{o} grado.
  \item \texttt{story4}: porcentaje que aprueba la prueba de lectura de 4.\textsuperscript{o} grado.
  \item \texttt{bs}: razón beneficios/salario, definida como \texttt{avgben/avgsal}.
 \texttt{lenrol}: logaritmo natural de \texttt{enrol}.
  \item \texttt{lstaff}: logaritmo natural de \texttt{staff}.
\end{itemize}

% Table created by stargazer v.5.2.3 by Marek Hlavac, Social Policy Institute. E-mail: marek.hlavac at gmail.com
% Date and time: Tue, Oct 14, 2025 - 11:46:19 AM
\begin{table}[!htbp] \centering 
  \caption{Tabla 1:Salarios de los profesores}  \label{tabla_reg}%\caption{Comparación: Regresiones} 

\begin{tabular}{@{\extracolsep{5pt}}lcc} 
\\[-1.8ex]\hline 
\hline \\[-1.8ex] 
 & \multicolumn{2}{c}{\textit{Dependent variable:}} \\ 
\cline{2-3} 
\\[-1.8ex] & \multicolumn{2}{c}{log(avgsal)} \\ 
\\[-1.8ex] & (1) & (2)\\ 
\hline \\[-1.8ex] 
 bs & $-$0.516$^{***}$ & $-$0.518$^{***}$ \\ 
  & (0.110) & (0.110) \\ 
  & & \\ 
 lenrol & $-$0.028$^{***}$ & $-$0.028$^{***}$ \\ 
  & (0.008) & (0.008) \\ 
  & & \\ 
 lstaff & $-$0.691$^{***}$ & $-$0.688$^{***}$ \\ 
  & (0.018) & (0.018) \\ 
  & & \\ 
 lunch & $-$0.001$^{***}$ & $-$0.001$^{***}$ \\ 
  & (0.0002) & (0.0002) \\ 
  & & \\ 
 math4 &  & 0.0004 \\ 
  &  & (0.0004) \\ 
  & & \\ 
 story4 &  & $$0.001 \\ 
  &  & (0.0005) \\ 
  & & \\ 
 Constant & 13.831$^{***}$ & 13.850$^{***}$ \\ 
  & (0.110) & (0.113) \\ 
  & & \\ 
\hline \\[-1.8ex] 
Observations & 1,848 & 1,848 \\ 
R$^{2}$ & 0.488 & 0.493 \\ 

\hline 
\hline \\[-1.8ex] 
\textit{Note:}  & \multicolumn{2}{r}{$^{*}$p$<$0.1; $^{**}$p$<$0.05; $^{***}$p$<$0.01} \\ 
\end{tabular} 
\end{table} 

Debiese incluir las variables \texttt{math4} y \texttt{story4} en la investigación? Sustente su respuesta mediante la formulación y explicación de un test de hipótesis que permita evaluar su relevancia estadística en el modelo ($
F_{0.05, (2, 1841)} \approx 3.00$).


{\color{blue} Se desea contrastar la hipótesis nula de que las dos variables eliminadas no aportan información estadísticamente significativa al modelo. Para ello, se emplea el test $F$ de restricciones lineales:

\[
H_0: \beta_{math4} =\beta_{story4} =0 
\]

\[
H_1: \text{Al menos uno es } \neq 0
\]

Los resultados estimados son los siguientes:

\[
R^2_{UR} = 0.493, \quad R^2_{R} = 0.488, \quad n = 1848, \quad k_{UR} = 6, \quad q = 2
\]

El estadístico de prueba se calcula como:

\[
F = \frac{(R^2_{UR} - R^2_{R}) / q}{(1 - R^2_{UR}) / (n - k_{UR} - 1)}
\]

Sustituyendo los valores:

\[
F = \frac{(0.493 - 0.488)/2}{(1 - 0.493)/(1848 - 6 - 1)} 
   = \frac{0.005 / 2}{0.507 / 1841}
   = \frac{0.0025}{0.0002754} \approx 9
\]

Por lo tanto, el estadístico $F$ sigue una distribución:

\[
F_{(q,\, n - k_{UR} - 1)} = F_{(2,\,1841)}
\]

El valor crítico al $5\%$ de significancia es aproximadamente:

\[
F_{0.05, (2, 1841)} \approx 3.00
\]

Dado que:

\[
F_{calc} = 9 > 3.00
\]

se \textbf{rechaza la hipótesis nula} $H_0$. 

\bigskip

\noindent
\textbf Existe evidencia estadísticamente significativa al 5\% de que las dos variables contribuyen de manera conjunta a explicar la variabilidad de la variable dependiente. Debiese incluir ambas variables en la investigación.
}

 
 \item \textbf{RdeA9.1} (10 puntos)  
{Considere el siguiente modelo del tiempo de entrega en función de la carga de trabajo:}

\[
\text{tiempo} = 12 - 0.7\,\text{carga} + 0.02\,\text{carga}^2
\]

\text{donde el tiempo se mide en horas y la carga representa el número de pedidos procesados por día.}


a) (5 puntos) ¿Cuál es la interpretación de $\beta_1$ y $\beta_2$?

\color{blue}
 $\boldsymbol{\beta_1}$: representa el efecto lineal inicial de la \textit{carga} sobre el \textit{tiempo de entrega}. 
 
  
$\boldsymbol{\beta_2}$: 
  implica que la pendiente de la relación entre \textit{tiempo} y \textit{carga}
aumenta en $2\beta_2 $ por cada unidad adicional de \textit{carga}.

Nos permite capturar los efectos no lineales de la carga sobre el tiempo de entrega, captura la segunda derivada de tiempo con respecto carga.  Si los valores poblacionales tienen los signos de los estimadores respectivos 

\color{black}
{b) (5 puntos) Determine el nivel de carga que minimiza el tiempo de entrega.}

\color{blue}
\medskip
Para encontrar el punto mínimo, derivamos la función respecto de \textit{carga}:

\[
\frac{d\,\text{tiempo}}{d\,\text{carga}} = -0.7 + 0.04\,\text{carga}
\]

Igualamos a cero para obtener el punto crítico:

\[
-0.7 + 0.04\,\text{carga} = 0
\quad \Rightarrow \quad
0.04\,\text{carga} = 0.7
\quad \Rightarrow \quad
\text{carga}^* = \frac{0.7}{0.04} = 17.5
\]

\textbf{Verificación del mínimo:}
\[
\frac{d^2\,\text{tiempo}}{d\,\text{carga}^2} = 0.04 > 0
\quad \Rightarrow \quad
\text{carga}^* \text{ corresponde a un mínimo.}
\]

\textbf{Conclusión:} 
El tiempo de entrega es mínimo cuando la carga de trabajo es de 
\[
\boxed{\text{carga}^* = 17.5\ \text{pedidos por día.}}
\]


\color{black}


 
 \item \textbf{RdeA9.1} (10 puntos) 

{Suponga un modelo aplicado a las ventas dado por:}

\[
\ln(\mathrm{ventas}) = \beta_0 + \beta_1\,\mathrm{publicidad} + \beta_2\,\mathrm{precio} + u
\]

\noindent
\textbf{Definiciones y unidades:}
\begin{itemize}
    \item \textbf{ventas}: número de unidades vendidas (logaritmo natural)
    \item \textbf{publicidad}: inversión en publicidad (miles de dólares)
    \item \textbf{precio}: precio del producto (dólares por unidad)
\end{itemize}

\begin{center}
\begin{tabular}{lcc}
\toprule
\textbf{Variable} & \textbf{Coeficiente} & \textbf{Error Estándar} \\
\midrule
Intercepto   & 2.00***  & 0.19 \\
publicidad   & 0.058**  & 0.021 \\
precio       & -0.016*  & 0.008 \\
\bottomrule
\end{tabular}

\vspace{0.5em}

$R^2 = 0.42$ \quad $F = 28.5$ (gl = 2, 92) \\[0.3em]
\textit{Nota:} *p $<$ 0.1; **p $<$ 0.05; ***p $<$ 0.01
\end{center}

 a) (5 puntos) 
Realice una prueba de significancia individual para el parámetro $\beta_{1}$ asociado a la variable \texttt{publicidad}, considerando un nivel de significancia del 10\%. Explique sus resultados.

\color{blue}
\medskip
\textit{Hipótesis:}\quad
H_0:\beta_1=0 \qquad \text{vs} \qquad H_1:\beta_1\neq 0.

\textit{Información disponible:} el coeficiente de \texttt{publicidad} aparece con ``**'', lo que indica un \emph{p}-valor \textbf{menor que 0.05} (prueba bilateral).

\textit{Decisión:} como $p<0.05<0.10$, se \textbf{rechaza $H_0$} al 10\%.

\textit{Conclusión:} el parámetro $\beta_1$ es \textbf{significativamente distinto de cero} al 10\% (de hecho, también al 5\%); por lo tanto, la publicidad tiene un efecto estadísticamente significativo sobre $\ln(\mathrm{ventas})$.
\bigskip


%Realice la inferencia para el parámetro $\beta_1$ (publicidad) al 10\% de significancia.
% Formule las hipótesis nula y alternativa, calcule el estadístico t correspondiente, indique el valor crítico
% y concluya si el parámetro es significativamente distinto de cero.
\color{black}
b) (5 puntos) Un investigador sostiene que el efecto del precio sobre las ventas es mayor que cero
al 1\% de significancia. ¿Está usted de acuerdo? Justifique su respuesta utilizando los resultados del modelo.

\color{blue}
En el cuadro, el coeficiente estimado para \texttt{precio} es \(\hat\beta_2=-0.016\) y aparece con una sola estrella ``*'', lo que indica \(p<0.10\) (prueba bilateral). Con esta información:

\begin{itemize}
  \item El \textbf{signo es negativo}, por lo que la evidencia no respalda \( \beta_2>0 \).
  \item La significancia reportada es a \(10\%\) (no a \(1\%\)); por tanto, \textbf{no} es significativo al \(1\%\).
  \item En un contraste unidireccional \(H_0:\beta_2= 0\) vs \(H_1:\beta_2>0\), un estimador negativo va en contra de \(H_1\).
\end{itemize}

\textbf{Conclusión:} Como bajo H1 el valor de $\beta_2$ es positivo solo es evidencia en contra de la nula la cola positiva. en este caso al tener $\beta_2$ un valor negativo el estadístico t de prueba también será negativo y no estará en la cola positiva.  Con los resultados disponibles, no se puede sostener que el efecto del \texttt{precio} sea mayor que cero al 1\%.
\color{black}

 
 \item \textbf{RdeA9.1} (10 puntos)  

Se plantea el modelo:

$$votosA = \beta_0 + \beta_1 \cdot ln(\text{gasto A}) + \beta_2 \cdot ln(\text{gasto B}) + \beta_3 \cdot forpart + u$$

\begin{table}[H]
\centering
\caption{Resultados de la estimación del modelo de regresión}
\label{tab:regresion}
\begin{tabular}{lcccc}
\toprule
\textbf{Variable} & \textbf{Coeficiente} & \textbf{Error estándar} & \textbf{t-valor} & \textbf{Pr(>|t|)} \\
\midrule
Intercepto & 45.07893 & 3.92631 & 11.48 & $<$2e-16$^{***}$ \\
$\ln(\text{gasto A})$ & 6.08332 & 0.38215 & 15.92 & $<$2e-16$^{***}$ \\
$\ln(\text{gasto B})$ & -6.61542 & 0.37882 & -17.46 & $<$2e-16$^{***}$ \\
\textit{forpart} & 0.15196 & 0.06202 & 2.45 & 0.0153$^{*}$ \\
\midrule
\multicolumn{5}{l}{\textit{Signif. codes:} $^{***}p<0.001$, $^{**}p<0.01$, $^{*}p<0.05$} \\
\midrule
\multicolumn{5}{l}{Error estándar residual: 7.712 (169 g.l.)} \\
\multicolumn{5}{l}{$R^2$ = 0.7926 \quad $R^2_{ajustado}$ = 0.7889} \\
\multicolumn{5}{l}{Estadístico F = 215.2 (3; 169) \quad p-valor $<$ 2.2e-16} \\
\bottomrule
\end{tabular}
\end{table}

Donde $ln(\text{gasto A})$ y $ln(\text{gasto B})$ son el logaritmo natural del gasto en campaña del candidato A y B, respectivamente, $forpart$ es una medida de la fortaleza del partido del candidato A (porcentaje de votos que ovtuvo el partido A en la elección presidencial más reciente), y votosA el porcentaje de votos recibidos por el candidato A.\\

Además, se tiene que la matriz de varianzas-covarianzas es:

\begin{table}[H]
    \centering
    \begin{tabular}{c|c|c|c|c}
 & (Intercept) & gasto A & gasto B	& forpart\\\hline
(Intercept)	& 15.4158713 & -0.394937598	& -0.874063434 & -0.176167017\\
gasto A	& -0.3949376 & 0.146038608 & -0.002681489 & -0.006546373\\
gasto B	& -0.8740634 & -0.002681489	& 0.143504821 & 0.003577341\\
forpart	& -0.1761670 & -0.006546373	& 0.003577341 & 0.003846245\\
    \end{tabular}
    %\caption{Caption}
    %\label{tab:placeholder}
\end{table}

Plantee el test de hipótesis de que un aumento de un 1\% en los gastos de A es compensado por un aumento de 1\% en los gastos de B, determine la veracidad de lo anterior al 5\% y explique.\\
\color{blue}
\textbf{Respuesta:} El test de hipótesis planteado es el siguiente: 

\begin{align*}
    H_0 &: \beta_1 = -\beta_2\\
    H_1 &: \beta_1 \neq -\beta_2
\end{align*}

\begin{align*}
    H_0 &: \beta_1 +\beta_2=0\\
H_1 &: \beta_1 +\beta_2 \neq 0 
\end{align*}
Luego, el cálculo de $t$ es:
$$t = \frac{\widehat{\beta}_1 + \widehat{\beta}_2}{\sqrt{V(\widehat{\beta}_1 + \widehat{\beta}_2)}}$$

Reemplazando y aplicando la propiedad de la varianza de una suma\footnote{Los valores finales están calculados con todos los decimales, puede haber diferencias por aproximación}:

$$t = \frac{6.083 - 6.615}{\sqrt{0.146 + 0.144 + 2\cdot (-0.003)}} \approx \frac{0.53324}{−0.53210 }= −0.9978$$

Luego, el $|t_{calc}| < 1.96$, por lo tanto, no existe evidencia suficiente para rechazar la hipótesis nula y, efectivamente, un aumento de 1\% en el gasto en campaña del candidato A se ve contrarrestado por un aumento de 1\% en gasto de campaña del candidato B.
    \end{enumerate}
\color{black}
\section{Comentes y preguntas cortas (60 puntos)}

Responda solo 2 de las siguientes preguntas.

\begin{enumerate}
    \item \textbf{RdeA9.1} (30 puntos) 
Considere el siguiente modelo de regresión lineal con los supuestos estándares:

\[
Y_i = \beta_0 + \beta_1 X_{1i} + \beta_2 X_{2i} + u_i
\]

Se dispone de la siguiente información:

\[
X'X = 
\begin{pmatrix}
99 & 0 & 0 \\
0 & 120 & 60 \\
0 & 60 & 180
\end{pmatrix},
\qquad
X'Y = 
\begin{pmatrix}
396 \\
72 \\
276
\end{pmatrix}
\]

\[
\sum_{i=1}^{N} Y_i^2 = 2034, 
\qquad 
\sum_{i=1}^{N} (Y_i - \bar{Y})^2 = 450
\]

\begin{enumerate}
    \item (3 puntos) ¿Cuál es el tamaño de la muestra?\\

    {\color{blue}
    $N = 99$, que corresponde al elemento $a_{11}$ de la matriz $X'X$.
    } (3 puntos, respuesta directa)\\[0.25em]
    
    \item (13 puntos) Calcule los estimadores de MCO del modelo anterior.\\
    \textbf{Hint:} Si \( Z = \begin{pmatrix} A &{\bf 0} \\ {\bf 0} & B \end{pmatrix} \) con \( A \) y \( B \) matrices cuadradas, entonces
    \[
    Z^{-1} = \begin{pmatrix} A^{-1} & {\bf 0} \\ {\bf 0} & B^{-1} \end{pmatrix}.
    \]

    {\color{blue}
    Los estimadores MCO están dados por
    \[
    \hat{\beta} = (X'X)^{-1} X'Y.
    \]
    Planteo \textbf{(2 puntos)}:
    \[
    \hat{\beta} 
    = 
    \begin{pmatrix}
    99 & 0 & 0\\
    0 & 120 & 60\\
    0 & 60 & 180
    \end{pmatrix}^{-1}
    \begin{pmatrix}
    396\\
    72\\
    276
    \end{pmatrix}.
    \]
    Inversa por bloques \textbf{(6 puntos)}. Para el escalar $99$, $(99)^{-1}=\tfrac{1}{99}$. Para el bloque\footnote{Acá todo está calculado con todos los decimales, por lo que los valores finales podrían diferir}
    \(
    B=
    \begin{pmatrix}
    120 & 60\\
    60 & 180
    \end{pmatrix}
    \),
    su inversa es
    \[
    B^{-1}
    = \frac{1}{\det(B)}
    \begin{pmatrix}
    180 & -60\\
    -60 & 120
    \end{pmatrix}
    = \frac{1}{18000}
    \begin{pmatrix}
    180 & -60\\
    -60 & 120
    \end{pmatrix}
    =
    \begin{pmatrix}
    0.01 & -0.003333333\\
    -0.003333333 & 0.006666667
    \end{pmatrix}.
    \]
    Por lo tanto
    \[
    (X'X)^{-1}
    =
    \begin{pmatrix}
    \tfrac{1}{99} & 0 & 0\\
    0 & 0.01 & -0.003333333\\
    0 & -0.003333333 & 0.006666667
    \end{pmatrix}.
    \]
    Ahora, multiplicando \textbf{(2 puntos para llegar al vector)}:
    \[
    \hat{\beta}
    =
    \begin{pmatrix}
    \tfrac{1}{99} & 0 & 0\\
    0 & 0.01 & -0.003333333\\
    0 & -0.003333333 & 0.006666667
    \end{pmatrix}
    \begin{pmatrix}
    396\\
    72\\
    276
    \end{pmatrix}
    =
    \begin{pmatrix}
    4\\
    -0.2\\
    1.6
    \end{pmatrix}.
    \]
    Conclusión \textbf{(1 punto por cada estimador indicado, 3 total)}:
    \[
    \boxed{\hat{\beta}_0 = 4,\quad \hat{\beta}_1 = -0.2,\quad \hat{\beta}_2 = 1.6.}
    \]
    }\\[0.25em]
    
    \item (4 puntos) Indique la ecuación del modelo predicho e interprete.\\

    {\color{blue}
    Ecuación predicha (2 puntos):
    \[
    \hat{Y} = 4 \;-\; 0.2\,X_1 \;+\; 1.6\,X_2.
    \]
    Interpretación (2 puntos): Manteniendo todo lo demás constante, un aumento de una unidad en $X_1$ reduce $Y$ en $0.2$ unidades; un aumento de una unidad en $X_2$ incrementa $Y$ en $1.6$ unidades.
    }\\[0.25em]
    
    \item (5 puntos) ¿Qué debería suceder con $R^2$ si agregamos más variables independientes a la medición?\\

    {\color{blue}
    \[
    R^2 \;=\; 1 - \frac{\sum_{i=1}^N \hat{u}_i^{\,2}}{\sum_{i=1}^N (Y_i-\bar{Y})^2}.
    \]
    Al agregar regresores (misma muestra y MCO), la suma de residuos al cuadrado no puede aumentar porque se relaja la restricción del problema de minimización. Por tanto, $R^2$ es \textbf{no decreciente} tiende a aumentar al agregar nuevos regresores. (El $R^2$ ajustado sí puede disminuir si los nuevos regresores no aportan suficiente poder explicativo.)
    }

    \item (5 puntos) ¿Qué valor alcanza el $R^2$?\footnote{Hint: Recuerde que $SSE = \sum_{i=1}^N (\hat{Y_i} - \bar{Y})^2$ que de manera matricial queda como $SSE=\hat{Y}^T Y - N \bar{Y^2}$}\\

    {\color{blue}
    Partimos de los resultados estimados:
    \[
    \hat{\beta} =
    \begin{pmatrix}
    4\\
    -0.2\\
    1.6
    \end{pmatrix}, \qquad
    X'Y =
    \begin{pmatrix}
    396\\
    72\\
    276
    \end{pmatrix}, \qquad
    N = 99, \qquad
    \sum (Y_i - \bar{Y})^2 = 450.
    \]
    
    \textbf{Cálculo de la media de $Y$:}
    \[
    \bar{Y} = \frac{\sum Y_i}{N} = \frac{(X'Y)_1}{N} = \frac{396}{99} = 4.
    \]
    
    \textbf{Cálculo de $SSE$:}
    \[
    SSE = \hat{Y}^T Y - N \bar{Y}^2 = \hat{\beta}^T (X'Y) - N \bar{Y}^2.
    \]
    \[
    \hat{\beta}^T (X'Y) = (4, -0.2, 1.6)
    \begin{pmatrix}
    396\\
    72\\
    276
    \end{pmatrix}
    = 4(396) + (-0.2)(72) + 1.6(276)
    = 1584 - 14.4 + 441.6 = 2011.2.
    \]
    \[
    N \bar{Y}^2 = 99 \times 4^2 = 99 \times 16 = 1584.
    \]
    \[
    SSE = 2011.2 - 1584 = 427.2.
    \]
    
    \textbf{Cálculo de $R^2$ (4 puntos):}
    \[
    R^2 = \frac{SSE}{SST} = \frac{427.2}{450} = 0.949333\ldots
    \]
    \[
    \boxed{R^2 \approx 0.9493.}
    \]
    
    \textbf{Interpretación (1 punto):} El modelo explica aproximadamente el $94.9\%$ de la variabilidad total de $Y$. 
}
\end{enumerate}

    
 \item \textbf{RdeA9.1} (30 puntos) 
 
 Considere una ecuación para explicar los sueldos de los directores generales o CEO en términos de las ventas anuales de la empresa, el rendimiento sobre capital (\emph{roe}, en forma de porcentaje), y el rendimiento de las acciones de la empresa (\emph{ros}, en forma de porcentaje):
\[
\log(\text{salary}) \;=\; \beta_0 \;+\; \beta_1 \log(\text{sales}) \;+\; \beta_2 \,\text{roe} \;+\; \beta_3 \,\text{ros} \;+\; u.
\]


 Con los datos de \texttt{CEOSAL1}, empleando MCO se obtuvo la ecuación siguiente:
\[
\widehat{\log(\text{salary})}
= 4.32 \;+\; 0.280\,\log(\text{sales}) \;+\; 0.0174\,\text{roe} \;+\; 0.0024\,\text{ros}
\]
\[
\qquad\qquad \qquad(0.32)\qquad \qquad(0.035)\qquad\qquad(0.041)\qquad\quad (0.0054)
\]
\[
n = 209,\qquad R^2 = 0.283.
\]

\begin{enumerate}
\item (6 puntos) Establezca la hipótesis nula de que, controlando \textit{sales} y \textit{roe}, \textit{ros} no tiene efecto en el sueldo de los CEO. Establezca la alternativa de que un mejor desempeño de las acciones de la empresa incrementa el sueldo de los CEO.  Realice la prueba al nivel de significancia de 10\%.

\item (6 puntos) Establezca la hipótesis nula de que, controlando \textit{sales} y \textit{roe}, \textit{ros} no tiene efecto en el sueldo de los CEO.  Establezca la alternativa de que el rendimiento sobre capital de la
empresa incrementa el sueldo de los CEO. Realice la prueba al nivel de significancia de 10\%.

\item (6 puntos) ¿Cuál es el porcentaje de aumento de \textit{salary} que se pronostica si \textit{ros} aumenta 50 puntos? ¿Tienen \textit{ros} un efecto práctico grande sobre \textit{salary}? ¿Es este efecto estadísticamente significativo?

\item (6 puntos)  Realice un {test} de {significancia conjunta} de \(\{\textit{roe},\textit{ros}\}\) condicionando por \(\log(\text{sales})\), para ello considere el siguiente modelo: 


\[
\widehat{\log(\text{salary})}
= 4.82 \;+\; 0.50\,\log(\text{sales}) 
\]
\[
\qquad\qquad \qquad(0.22)\qquad \qquad(0.015)
\]
\[
n = 209,\qquad R^2 = 0.279.
\]


 ($
F_{0.90}(2,205)\approx 2.35,
\quad
F_{0.95}(2,205)\approx 3.04.
$)

\item (6 puntos)  ¿Incluiría usted \textit{roe} y \textit{ros} en el modelo final que explica las compensaciones de los CEO en términos del desempeño de la empresa? Explique.
\end{enumerate}

{\color{blue} \begin{enumerate}
%------------------------------------------------
\item \textbf{Prueba individual para \texttt{ros}}

\[
H_0: \beta_3 = 0 \qquad \text{vs.} \qquad H_1: \beta_3 > 0
\]

El estadístico de prueba es:
\[
t = \frac{\hat{\beta}_3}{\text{se}(\hat{\beta}_3)} = \frac{0.0024}{0.0054} = 0.44
\]

  
 Bastaba con calcular el estadístico y  analizar intuitivamente el valor obtenido es muy bajo por lo tanto es evidencia a favor de la nula, no se entrega el valor crítico al 10\%, pero si para el 5\%.  
 
Con $n=209$ observaciones y $k=3$ regresores, los grados de libertad son $205$.
El valor crítico al 10\% (una cola) es $t_{0.90,205} \approx 1.29$.
Con $t=0.44 < 1.29$, \textbf{no se rechaza $H_0$}.  


%------------------------------------------------
\item \textbf{Prueba individual para \texttt{roe}}

\[
H_0: \beta_2 = 0 \qquad \text{vs.} \qquad H_1: \beta_2 > 0
\]

\[
t = \frac{\hat{\beta}_2}{\text{se}(\hat{\beta}_2)} = \frac{0.0174}{0.041} = 0.42
\]

Bastaba con calcular el estadístico y  analizar intuitivamente el valor obtenido es muy bajo por lo tanto es evidencia a favor de la nula, no se entrega el valor crítico al 10\%, pero si para el 5\%.  

 $t = 0.42 < 1.29$, \textbf{no se rechaza $H_0$}.  
No hay evidencia estadística al 10\% de que un mayor rendimiento sobre el capital (\texttt{roe}) aumente el salario de los CEO.

%------------------------------------------------
\item \textbf{Efecto de un aumento de 50 puntos en \texttt{ros}}

En un modelo semilogarítmico:
\[
\Delta \log(\text{salary}) \approx \hat{\beta}_3 \cdot \Delta \texttt{ros}
\]
\[
\Rightarrow 100 \times 0.0024 \times 50 = 12.0\% \text{ (aproximación lineal)}
\]
\[
\text{Ajuste exacto: } 100 \times [\exp(0.0024 \times 50) - 1] \approx 12.75\%
\]

El efecto estimado no es estadísticamente significativo (\(t = 0.44\)), por lo que no puede concluirse que un aumento en \texttt{ros} eleve el salario de los CEO.

%------------------------------------------------
\item \textbf{Prueba de significancia conjunta para \texttt{roe} y \texttt{ros}}

\[
H_0: \beta_2 = \beta_3 = 0 \qquad \text{vs.} \qquad H_1: \text{al menos uno distinto de cero.}
\]

Se tiene:
\[
R^2_{UR} = 0.283, \quad R^2_{R} = 0.279, \quad q = 2, \quad n = 209, \quad k_{UR}=3.
\]

El estadístico \(F\) es:
\[
F = \frac{(R^2_{UR}-R^2_R)/q}{(1-R^2_{UR})/(n-k_{UR}-1)} 
  = \frac{(0.283-0.279)/2}{(1-0.283)/205} 
  = \frac{0.002}{0.00350} = 0.57.
\]

Valores críticos: \(F_{0.90}(2,205) = 2.35\)  

\item Ni individualmente ni conjuntamente las variables \texttt{roe} y \texttt{ros} resultan estadísticamente significativas al 10\%.  
El incremento en el coeficiente de determinación es mínimo ($R^2$: 0.279 $\rightarrow$ 0.283). 
Por lo tanto, desde un criterio de parsimonia y significancia, \textbf{no se recomienda incluir \texttt{roe} y \texttt{ros}} en el modelo final que explica el salario de los CEO en función del desempeño de la empresa.  


\end{enumerate}
}



 \item \textbf{RdeA9.1} (30 puntos)  

 Considere el siguiente modelo estimado a partir de una muestra de 120 individuos en Chile.  
El objetivo es analizar cómo la educación y el ingreso influyen en el estado de salud auto-reportado (escala 1–10):

\[
\mathrm{salud}_i = \beta_0 + \beta_1\,\mathrm{educ}_i + \beta_2\,\ln(\mathrm{ingreso}_i) + \varepsilon_i
\]

\noindent
\textbf{Descripción de variables:}
\begin{itemize}
    \item \textbf{salud}: Puntaje de salud auto-reportado (1 = muy mala, 10 = excelente).
    \item \textbf{educ}: Años de educación formal del individuo.
    \item \textbf{$\ln(\mathrm{ingreso})$}: Logaritmo natural del ingreso mensual del individuo (ingreso medido en miles de CLP).
\end{itemize}

\noindent
\textbf{Tabla 1: Resultados de la regresión}

\begin{table}[h!]
\centering
\begin{tabular}{lcc}
\toprule
\textbf{Variable} & \textbf{Coeficiente} & \ \\
\midrule
Intercepto         & 3.90***   \\
educación          & 0.120**   \\
$\ln(\mathrm{ingreso})$ & 0.350**  \\
\bottomrule
\end{tabular}
\vspace{0.5em}

\noindent
$R^2 = 0.36$ \\[0.3em]

\end{table}

\noindent
La matriz var–covar estimada de los coeficientes $\hat{\beta}$ es la siguiente:

\[
\widehat{\mathrm{Var}}(\hat{\beta}) =
\begin{bmatrix}
0.1600 & -0.0020 & 0.0010 \\
-0.0020 & 0.0025 & -0.0003 \\
0.0010 & -0.0003 & 0.0196
\end{bmatrix}
\]



---

\begin{flushleft}
\textbf a) (10 puntos) Interprete los signos y la significancia estadística de los coeficientes estimados para \textbf{educación} e \textbf{ingreso}.  
Explique si los resultados son consistentes con lo que se esperaría desde un punto de vista económico o social.

\color{blue}
\bigskip


\emph{Educación.} $\widehat\beta_1=0.120>0$: un año adicional de educación se asocia, en promedio, a un aumento de $0.120$ puntos en la salud (escala 1–10), ceteris paribus. 
\[
t=\frac{0.120}{0.05}=2.40 \quad(\text{dos colas, }df=117)\ \Rightarrow\ \text{significativo al 5\% (coincide con ``**'')}.
\]

\emph{ln(Ingreso)} $\widehat\beta_2=0.350>0$: como el regresor es $\ln(\text{ingreso})$, una variación de $1\%$ en el ingreso ($\Delta\ln=0.01$) se asocia a $0.0035$ puntos más de salud ($0.350\times0.01$), manteniendo lo demás constante.
\[
t=\frac{0.350}{0.14}=2.50 \quad(\text{dos colas, }df=117)\ \Rightarrow\ \text{significativo al 5\% (``**'')}.
\]

\emph{Coherencia.} Ambos signos positivos son consistentes con la expectativa económica/social: más educación y mayor ingreso suelen relacionarse con mejor salud auto-reportada.


\color{black}


\vspace{0.5em}
\textbf b) (10 puntos) Testee la siguiente hipótesis: el efecto del ingreso sobre la salud es el triple del efecto de la educación sobre la salud, frente a la alternativa de que es diferente. Suponga $\alpha = 5\%$.

\color{blue}

Lo analizamos con el valor obtenido para el logaritmo del ingreso.

\textbf{(b) Prueba de hipótesis: } 
$H_0: \beta_2 = 3\beta_1$ 
\quad vs \quad 
$H_1: \beta_2 \neq 3\beta_1$

$H_0: \beta_2 - 3\beta_1=0$ 
\quad vs \quad 
$H_1: \beta_2 - 3\beta_1\neq 0$
\medskip
El estadístico es:
\[
t = \frac{(\hat{\beta}_2 - 3\hat{\beta}_1) - 0}
{\sqrt{\mathrm{Var}(\hat{\beta}_2 - 3\hat{\beta}_1)}}.
\]

\textbf{Cálculos:}

\[
\hat{\beta}_2 - 3\hat{\beta}_1 = 0.350 - 3(0.120) = -0.010.
\]

\[
\mathrm{Var}(\hat{\beta}_2 - 3\hat{\beta}_1)
= \mathrm{Var}(\hat{\beta}_2)
+ 9\,\mathrm{Var}(\hat{\beta}_1)
- 6\,\mathrm{Cov}(\hat{\beta}_1,\hat{\beta}_2)
= 0.0196 + 9(0.0025) - 6(-0.0003) = 0.0439.
\]

\[
s.e.(\hat{\beta}_2 - 3\hat{\beta}_1)
= \sqrt{0.0439} = 0.2096.
\]

\[
t = \frac{-0.010}{0.2096} = -0.048.
\]

Como $|t| < t_{0.975,117} \approx 1.96$, \textbf{no se rechaza $H_0$ al 5\%}.
\color{black}

\vspace{0.5em}
\textbf c) (10 puntos) Pruebe conjuntamente la hipótesis de que la educación y el ingreso no influyen en el estado de salud. Suponga $\alpha = 10\%$.($F_{0.90}(2,117) \approx 2.30$)
\color{blue}

\textbf{Prueba conjunta (educación e ingreso no influyen en salud), $\alpha=10\%$}

\medskip
Hipótesis:
\[
H_0:\ \beta_1=\beta_2=0
\qquad \text{vs} \qquad
H_1:\ \text{al menos uno}\neq 0.
\]

Usamos el estadístico $F$ global con $R^2=0.36$, $n=120$, $k=2$:
\[
F=\frac{(R^2/k)}{(1-R^2)/(n-k-1)}
=\frac{0.36/2}{0.64/117}
=\approx 32.9,
\qquad F \sim F(2,117).
\]

Valor crítico (bilateral en términos del $F$) a $\alpha=0.10$:
\[
F_{0.90}(2,117)\approx 2.30.
\]

Como $32.9 \gg 2.30$, \textbf{rechazamos $H_0$} al 10\%.

\[
\textbf{Conclusión:} \ \text{en conjunto, educación e ingreso sí influyen en el estado de salud.}




\end{flushleft}
\color{black}
\section{Formulario}
\begin{itemize}

\item  
{\color{red}
$$\hat{\beta}_0 = \bar{y} - \hat{\beta}_1 \bar{x} $$
$$\hat{\beta}_1 = \frac{\sum_{i=1}^{n}(x_i - \bar{x})(y_i-\bar{y})}{\sum_{i=1}^{n}(x_i - \bar{x})^2}$$}

 {\color{red} \[ \hat{\beta} = (X^{\top}X)^{-1}X^{\top}y \]}

\item 
	{\color{red}	\[ SST = \sum_{i=1}^{n} (y_i - \bar{y})^2 \quad SSE = \sum_{i=1}^{n} (\hat{y}_i - \bar{y})^2  \quad SSR = \sum_{i=1}^{n} \hat{u}_i^2\]}
		

    \item {\color{red} \[ R^2  = 1 - \frac{SSR}{SST} \quad \quad \bar{R}^2 =  1 - \frac{(n-1)SSR}{(n-k-1)SST} \]}

 \item {\color{red} \[ t_{\hat{\beta_j}}=\frac{\hat{\beta}_j - \beta_j}{\text{s.e.}(\hat{\beta}_j)} \sim t_{n-k-1} \]}

 \item {\color{red}\[
    F = \frac{SSR_{r} - SSR_{ur}}{SSR_{ur}} \cdot \frac{n - k - 1}{j} \sim F(j, n - k - 1)
    \]}
\end{itemize}   

\begin{figure}[H]
    \centering
    \includegraphics[width=1.05\linewidth]{certamenes/norm_bilat_ren.png}
\end{figure}
 
    \end{enumerate}
    \end{document}